
\documentclass[11pt]{article}
\usepackage{amsmath,amssymb,amsthm}
\usepackage[a4paper,margin=1in]{geometry}

\newtheorem{lemma}{Lemma}
\newtheorem{corollary}{Corollary}

\begin{document}

\section*{Horizontal Hessian Bound and the Gribov Region}

Recall that in exponential coordinates $U_b = \exp(i a g A_b)$ near the identity, the effective action
\begin{equation*}
S_{\mathrm{eff}}(U) = \beta S_W(U) + S_{\mathrm{Haar}}(U)
\end{equation*}
has a horizontal Hessian (restricted to the subspace orthogonal to gauge orbits)
\begin{equation*}
H_H(U) := \mathrm{Hess}\, S_{\mathrm{eff}}(U)\big|_{T^{\mathrm{hor}}_U \mathcal{C}}
\end{equation*}
which admits the lower bound
\begin{equation}
\label{eq:HH-lower-bound}
\langle A, H_H(U) A \rangle \;\ge\; \bigl(c_0 a^2 g^2 - \beta C_V\bigr)\,\|A\|^2
\qquad \text{for all } A \in T^{\mathrm{hor}}_U \mathcal{C}.
\end{equation}
Here $c_0 > 0$ is the Haar mass coefficient, $a$ the lattice spacing, $g$ the bare coupling, $\beta$ the plaquette coupling, and $C_V \ge 0$ a uniform bound on the potential term from the Wilson action.

Define the smallest horizontal eigenvalue of $H_H(U)$ by
\begin{equation*}
\lambda_{\min}(U) \;=\; \min_{A \in T^{\mathrm{hor}}_U \mathcal{C},\, A \neq 0}
\frac{\langle A, H_H(U) A \rangle}{\|A\|^2}\,.
\end{equation*}

\begin{lemma}[Rayleigh-quotient bound for the horizontal Hessian]
\label{lem:lambda-min-bound}
If \eqref{eq:HH-lower-bound} holds for all horizontal $A$, then
\begin{equation}
\label{eq:lambda-min-bound}
\lambda_{\min}(U) \;\ge\; c_0 a^2 g^2 - \beta C_V.
\end{equation}
\end{lemma}

\begin{proof}
By definition of $\lambda_{\min}(U)$ and the Rayleigh quotient,
\begin{equation*}
\lambda_{\min}(U)
= \min_{A \in T^{\mathrm{hor}}_U \mathcal{C},\, A \neq 0}
\frac{\langle A, H_H(U) A \rangle}{\|A\|^2}.
\end{equation*}
The assumed inequality \eqref{eq:HH-lower-bound} gives, for every nonzero horizontal $A$,
\begin{equation*}
\frac{\langle A, H_H(U) A \rangle}{\|A\|^2}
\;\ge\; c_0 a^2 g^2 - \beta C_V.
\end{equation*}
Taking the minimum over all admissible $A$ cannot decrease the right-hand side, hence
\begin{equation*}
\lambda_{\min}(U) \;\ge\; c_0 a^2 g^2 - \beta C_V,
\end{equation*}
which is \eqref{eq:lambda-min-bound}.
\end{proof}

This immediately yields a clean condition for strict positivity.

\begin{corollary}[Strict convexity and Gribov region]
\label{cor:Gribov-region}
Assume $c_0 a^2 g^2 - \beta C_V > 0$. Then for every configuration $U$ for which
\eqref{eq:HH-lower-bound} holds, the horizontal Hessian $H_H(U)$ is strictly positive definite, and
\begin{equation*}
\lambda_{\min}(U) \;\ge\; \rho_\ast(a,g,\beta) := c_0 a^2 g^2 - \beta C_V \;>\; 0.
\end{equation*}
The set of such configurations,
\begin{equation*}
\Omega_G \;=\; \bigl\{\,U \in \mathcal{C} \;:\; \lambda_{\min}(U) > 0\,\bigr\},
\end{equation*}
is (in the usual gauge-fixing picture) identified with the Gribov region, while the hypersurface
\begin{equation*}
\partial\Omega_G \;=\; \{\, U : \lambda_{\min}(U) = 0 \,\}
\end{equation*}
is the Gribov horizon.
\end{corollary}

\begin{remark}
If $c_0 a^2 g^2 - \beta C_V < 0$, the bound \eqref{eq:lambda-min-bound} becomes non-informative: it allows the possibility of negative eigenvalues but does not force them. In that regime our present estimate cannot certify stability, and one expects either the appearance of zero-modes (on the horizon) or actual negative modes beyond it, depending on the detailed structure of $H_H(U)$.

In the finite-cutoff mass-gap argument, we work in the parameter regime where
\begin{equation*}
c_0 a^2 g^2 \;>\; \beta C_V,
\end{equation*}
so that a uniform strictly positive lower bound $\rho_\ast(a,g,\beta)$ exists on the horizontal Hessian. This is exactly the input needed for the Bakry--\'Emery curvature condition and the associated functional inequalities.
\end{remark}

\end{document}
